% Generated draft: requires mathematical and source review.
\documentclass[11pt]{article}
\usepackage[margin=25mm]{geometry}
\usepackage{fontspec}
\setmainfont{Latin Modern Roman}
\usepackage{amsmath,amssymb,mathtools,unicode-math}
\setmathfont{Latin Modern Math}
\usepackage{enumitem,needspace,xurl,hyperref}
\setlength{\parindent}{0pt}
\setlength{\parskip}{5pt}
\setlength{\emergencystretch}{3em}
\newcommand{\N}{\mathbb{N}}
\newcommand{\Z}{\mathbb{Z}}
\newcommand{\Q}{\mathbb{Q}}
\newcommand{\R}{\mathbb{R}}
\newcommand{\C}{\mathbb{C}}
\newcommand{\F}{\mathbb{F}}
\newcommand{\A}{\mathbb{A}}
\newcommand{\PP}{\mathbb{P}}
\newcommand{\E}{\mathbb{E}}
\newcommand{\eps}{\varepsilon}
\newcommand{\dd}{\,\mathrm{d}}
\begin{document}
\section*{UnsolvedMath ID: 701781}
\textbf{Source problem code:} catalogue entry
\subsection{Definitions and mathematical statement}
Let \(\pi\) be any set of primes. A finite \(\pi\)-group is a finite group whose order has no prime divisors outside \(\pi\). A \(\pi\)-subgroup of a finite group \(G\) is a subgroup that is a \(\pi\)-group; it is maximal if it is maximal under inclusion among the \(\pi\)-subgroups of \(G\). Following the supplied source, call \(G\) a \(D_{\pi}\)-group if any two maximal \(\pi\)-subgroups of \(G\) are conjugate in \(G\). A Hall \(\pi\)-subgroup \(H\leq G\) is a \(\pi\)-subgroup such that no prime in \(\pi\) divides \([G:H]\). Write \(O_{\pi}(G)\) for the largest normal \(\pi\)-subgroup of \(G\), and use the convention \(H^x=x^{-1}Hx\). The soluble radical \(R(G)\) is the largest soluble normal subgroup of \(G\). A composition factor means a factor in a composition series of \(G\).

\begin{enumerate}
\item[(a)] Is it true that, for every set of primes \(\pi\), every finite \(D_{\pi}\)-group \(G\), and every Hall \(\pi\)-subgroup \(H\leq G\), there exist \(x,y,z\in G\) such that
\[
O_{\pi}(G)=H\cap H^x\cap H^y\cap H^z?
\]
\item[(b)] Is it true that, for every set of primes \(\pi\), every finite \(D_{\pi}\)-group \(G\) whose nonabelian simple composition factors are all sporadic simple groups or alternating groups \(A_n\) with \(n\geq 5\), and every Hall \(\pi\)-subgroup \(H\leq G\), there exist \(x,y\in G\) such that
\[
H\cap H^x\cap H^y=O_{\pi}(G)?
\]
\item[(c)] Is it true that, for every set of primes \(\pi\), every finite \(D_{\pi}\)-group \(G\) satisfying \(R(G)=1\) and having all nonabelian simple composition factors sporadic, and every Hall \(\pi\)-subgroup \(H\leq G\), there exists \(g\in G\) such that
\[
H\cap H^g=O_{\pi}(G)?
\]
\end{enumerate}

The supplied record labels the overall problem partially solved and reports affirmative solutions to parts (b) and (c). It supplies no solution notice for part (a). These are supplied source-status qualifications, not independently verified mathematical or literature assessments.
\subsection{Short English statement}
The problem asks whether the largest normal \(\pi\)-subgroup is the intersection of four conjugates of any given Hall \(\pi\)-subgroup in a finite \(D_{\pi}\)-group. Under the composition-factor restrictions in (b) and (c), it asks whether three or two conjugates, respectively, suffice. The supplied source reports affirmative solutions to (b) and (c).
\subsection{Sources}
\begin{itemize}
\item Local UnsolvedMath registry: permanent ID \(701781\), title \emph{Kourovka Notebook Problem 20.123}, exact source code \texttt{catalogue entry}. The supplied registry status is \texttt{partially\_solved}.
\item V. I. Zenkov, proposer, 2022. \emph{The Kourovka Notebook (Unsolved Problems in Group Theory)}, 20th issue (2022), Problem 20.123; also the supplied 21st issue edition (2026), \texttt{arXiv:1401.0300v48}. The supplied collection description names E. I. Khukhro and V. D. Mazurov as editors.
\item Part (b): the supplied source reports an affirmative solution by I. N. Belousov and V. I. Zenkov, preprint, 2025, in Russian: \url{https://kourovkanotebookorg.wordpress.com/wp-content/uploads/2025/09/20.123b.pdf}.
\item Part (c): the supplied source reports an affirmative solution by I. N. Belousov and V. I. Zenkov, \emph{Trudy Inst. Mat. Mekh. Ural. Otdel. Ross. Akad. Nauk}, \textbf{31}, no. 1 (2025), 19--35, in Russian. No article title or link is supplied for this reference.
\item The supplied status review is dated \texttt{2026-10-16}, identifies its evidence as \texttt{maintained\_source}, and attributes the solution notices to the 21st issue (October 2026), \texttt{arXiv:1401.0300v48}. This is supplied dated curation, not an independently verified assessment. Its quotation of the part (c) reference is truncated; the complete journal details above come from the supplied background field.
\end{itemize}
\end{document}
